% Options for packages loaded elsewhere
\PassOptionsToPackage{unicode}{hyperref}
\PassOptionsToPackage{hyphens}{url}
\PassOptionsToPackage{dvipsnames,svgnames,x11names}{xcolor}
%
\documentclass[
  13pt,
  a4paperpaper,
  12pt]{article}
\usepackage{amsmath,amssymb}
\usepackage{setspace}
\usepackage{iftex}
\ifPDFTeX
  \usepackage[T1]{fontenc}
  \usepackage[utf8]{inputenc}
  \usepackage{textcomp} % provide euro and other symbols
\else % if luatex or xetex
  \usepackage{unicode-math} % this also loads fontspec
  \defaultfontfeatures{Scale=MatchLowercase}
  \defaultfontfeatures[\rmfamily]{Ligatures=TeX,Scale=1}
\fi
\usepackage{lmodern}
\ifPDFTeX\else
  % xetex/luatex font selection
  \setmainfont[]{Times New Roman}
  \setmonofont[]{Courier New}
\fi
% Use upquote if available, for straight quotes in verbatim environments
\IfFileExists{upquote.sty}{\usepackage{upquote}}{}
\IfFileExists{microtype.sty}{% use microtype if available
  \usepackage[]{microtype}
  \UseMicrotypeSet[protrusion]{basicmath} % disable protrusion for tt fonts
}{}
\makeatletter
\@ifundefined{KOMAClassName}{% if non-KOMA class
  \IfFileExists{parskip.sty}{%
    \usepackage{parskip}
  }{% else
    \setlength{\parindent}{0pt}
    \setlength{\parskip}{6pt plus 2pt minus 1pt}}
}{% if KOMA class
  \KOMAoptions{parskip=half}}
\makeatother
\usepackage{xcolor}
\usepackage[margin=1.5cm,top=1.5cm,bottom=1.5cm,left=1.5cm,right=1.5cm,includeheadfoot]{geometry}
\usepackage{listings}
\newcommand{\passthrough}[1]{#1}
\lstset{defaultdialect=[5.3]Lua}
\lstset{defaultdialect=[x86masm]Assembler}
\setlength{\emergencystretch}{3em} % prevent overfull lines
\providecommand{\tightlist}{%
  \setlength{\itemsep}{0pt}\setlength{\parskip}{0pt}}
\setcounter{secnumdepth}{3}
% LaTeX Preamble for Insect Perception Research
\documentclass[12pt,a4paper]{article}
\usepackage[utf8]{inputenc}
\usepackage[T1]{fontenc}
\usepackage[margin=1in, top=1in, bottom=1in]{geometry}
\usepackage{graphicx}
\usepackage{amsmath}
\usepackage{amsfonts}
\usepackage{amssymb}
\usepackage{booktabs}
\usepackage{xcolor}
\usepackage{listings}
\usepackage{fancyvrb}
\usepackage{setspace}
\usepackage{ragged2e}
\usepackage{etoolbox}
\usepackage{chngcntr}
\usepackage[font=small,labelfont=bf,textfont=it]{caption}
\usepackage{titlesec}
\usepackage{amsthm}
\usepackage[nameinlink,capitalise]{cleveref}

% Hyperref package - properly configured for cross-references
\usepackage[
    colorlinks=true,
    linkcolor=blue,
    citecolor=blue,
    urlcolor=blue,
    linkbordercolor=blue,
    urlbordercolor=blue,
    citebordercolor=blue
]{hyperref}

% Custom colors for code listings
\definecolor{codebg}{RGB}{245, 245, 245}
\definecolor{codeborder}{RGB}{200, 200, 200}
\definecolor{codefg}{RGB}{50, 50, 50}

% Code listing setup
\lstset{
    backgroundcolor=\color{codebg},
    basicstyle=\normalsize\ttfamily\color{codefg},
    breakatwhitespace=false,
    breaklines=true,
    captionpos=b,
    frame=single,
    framerule=0.5pt,
    framesep=5pt,
    rulecolor=\color{codeborder},
    numbers=left,
    numbersep=5pt,
    numberstyle=\small\color{codefg}
}

% Font and spacing for accessibility
% Using standard LaTeX fonts for better compatibility
\renewcommand{\familydefault}{\rmdefault}
\renewcommand{\sfdefault}{\rmdefault}

% Single spacing for academic format
\setstretch{1.0}
\setlength{\parindent}{0.5in}
\setlength{\parskip}{0em}

% Left-justify the body text globally
\RaggedRight
\setlength{\RaggedRightParindent}{0.5in}

% Ensure body text is fully justified and consistent
\setlength{\parindent}{0.5in}

% Comprehensive numbering for figures, tables, and equations
\counterwithout{figure}{section}
\counterwithout{table}{section}
\counterwithout{equation}{section}

% Figure and table caption formatting
\captionsetup{justification=raggedright,singlelinecheck=false}
\captionsetup[figure]{position=bottom,skip=10pt}
\captionsetup[table]{position=top,skip=10pt}
% Left-align content inside all figure environments
\AtBeginEnvironment{figure}{\raggedright}

% Configure cleveref names for better cross-references
\crefname{figure}{Figure}{Figures}
\crefname{table}{Table}{Tables}
\crefname{equation}{Equation}{Equations}
\crefname{section}{Section}{Sections}
\crefname{chapter}{Chapter}{Chapters}

% Equation numbering and formatting
\newtheorem{theorem}{Theorem}[section]
\newtheorem{lemma}[theorem]{Lemma}
\newtheorem{proposition}[theorem]{Proposition}
\newtheorem{corollary}[theorem]{Corollary}
\newtheorem{definition}[theorem]{Definition}
\newtheorem{example}[theorem]{Example}
\newtheorem{remark}[theorem]{Remark}

% Section formatting for academic style
\titlespacing{\section}{0pt}{12pt}{6pt}
\titlespacing{\subsection}{0pt}{10pt}{4pt}
\titlespacing{\subsubsection}{0pt}{8pt}{3pt}

% Title formats for consistent academic appearance
\titleformat{\section}{\large\bfseries}{\thesection}{1em}{}
\titleformat{\subsection}{\normalsize\bfseries}{\thesubsection}{1em}{}
\titleformat{\subsubsection}{\normalsize\itshape}{\thesubsubsection}{1em}{}

% Title and author information
\title{When do bugs see (infra)red? On the Visual and Infra-red in the Insect Perceptual Apparatus}
\author{Tucker C. Chambers\thanks{Corresponding author} \\
\small Email: tucker.chambers@research.edu \\
\small ORCID: 0000-0001-6232-9096
\and
Daniel A. Friedman \\
\small Email: daniel@activeinference.institute \\
\small ORCID: 0000-0001-6232-9096}
\ifLuaTeX
  \usepackage{selnolig}  % disable illegal ligatures
\fi
\IfFileExists{bookmark.sty}{\usepackage{bookmark}}{\usepackage{hyperref}}
\IfFileExists{xurl.sty}{\usepackage{xurl}}{} % add URL line breaks if available
\urlstyle{same}
\hypersetup{
  colorlinks=true,
  linkcolor={blue},
  filecolor={blue},
  citecolor={blue},
  urlcolor={blue},
  pdfcreator={LaTeX via pandoc}}

\author{}
\date{}

\begin{document}

{
\hypersetup{linkcolor=black}
\setcounter{tocdepth}{3}
\tableofcontents
}
\setstretch{1.0}
\section{Appendix A: Sensilla Array Directionality and Beam
Patterns}\label{sec:app_sensilla_array}

\subsection{Objective}\label{objective}

Electromagnetic antenna modeling for sensilla arrays:
circular/log‑periodic designs, element patterns, mutual coupling, 2D
radiation patterns, morphology analysis, and frequency‑response
characterization for directional olfactory detection.

\subsection{Methods (src)}\label{methods-src}

\begin{itemize}
\tightlist
\item
  \passthrough{\lstinline!src/case\_studies/sensilla\_array\_directionality.py!}

  \begin{itemize}
  \tightlist
  \item
    \passthrough{\lstinline!design\_circular\_array(n\_elements: int, radius\_m: float, wavelength\_m: float) -> np.ndarray!}
  \item
    \passthrough{\lstinline!sensilla\_element\_pattern(sensilla\_type: str, frequency\_hz: float, dimensions: dict) -> np.ndarray!}
  \item
    \passthrough{\lstinline!mutual\_coupling\_matrix(positions: np.ndarray, wavelength\_m: float) -> np.ndarray!}
  \item
    \passthrough{\lstinline!array\_pattern\_2d(positions: np.ndarray, element\_patterns: np.ndarray, frequency: float, coupling: np.ndarray) -> np.ndarray!}
  \item
    \passthrough{\lstinline!analyze\_sensilla\_morphology(dimensions: np.ndarray, frequency\_range: np.ndarray) -> dict!}
  \item
    \passthrough{\lstinline!frequency\_response\_analysis(array\_config: dict, freq\_range: np.ndarray) -> dict!}
  \item
    \passthrough{\lstinline!compute\_beam\_pattern(wavelengths: np.ndarray, positions: np.ndarray, gains: np.ndarray) -> np.ndarray!}
  \item
    \passthrough{\lstinline!array\_gain(pattern: np.ndarray) -> float!}
  \item
    \passthrough{\lstinline!design\_log\_periodic\_array(min\_len: float, max\_len: float, tau: float, count: int) -> np.ndarray!}
  \end{itemize}
\end{itemize}

\subsection{Script and outputs}\label{script-and-outputs}

\begin{itemize}
\tightlist
\item
  Script:
  \passthrough{\lstinline!scripts/generate\_sensilla\_array\_directionality.py!}
\item
  Data:
  \passthrough{\lstinline!output/data/sensilla\_array\_comprehensive.npz!}
\item
  Figure:
  \passthrough{\lstinline!output/figures/sensilla\_array\_comprehensive\_analysis.png!}
\item
  Caption metadata:
  \passthrough{\lstinline!output/figures/sensilla\_array\_comprehensive\_analysis.caption.txt!}
\end{itemize}

\subsection{Figure}\label{figure}

\begin{figure}[h]
\centering
\includegraphics[width=0.9\textwidth]{../output/figures/sensilla_array_comprehensive_analysis.png}
\caption{Deterministic sensilla array analysis produced by `scripts/generate_sensilla_array_directionality.py` (seed=42): circular, linear, and log‑periodic array configurations; 2D radiation patterns; element patterns (dipole, monopole, patch); mutual coupling effects; morphology‑to‑resonance comparisons; and frequency‑response characterization.}
\label{fig:app_sensilla_beam}
\end{figure}

\subsection{Equation references}\label{equation-references}

\begin{itemize}
\tightlist
\item
  Effective aperture: see \eqref{eq:effective_aperture}
\item
  Gain pattern: see \eqref{eq:gain_pattern}
\end{itemize}

\subsection{Reproducibility}\label{reproducibility}

\begin{enumerate}
\def\labelenumi{\arabic{enumi}.}
\tightlist
\item
  Run:
  \passthrough{\lstinline!python3 scripts/generate\_sensilla\_array\_directionality.py!}
\item
  Artifacts: \passthrough{\lstinline!output/data/!} and
  \passthrough{\lstinline!output/figures/!}
\item
  Deterministic seed:
  \passthrough{\lstinline!src/config.set\_random\_seed(42)!}
\end{enumerate}

\subsection{Cross‑references}\label{crossreferences}

\begin{itemize}
\tightlist
\item
  Methods: \cref{sec:methodology}
\item
  Symbols: \cref{sec:symbols_glossary}
\item
  Math: \cref{sec:mathematical_appendix}
\end{itemize}

\end{document}
